% 迁移来源：NSFC_2026_CIESM-STSR/extraTex/2.3.年度研究计划.tex
\subsubsection{年度研究计划}

研究期限为\CiesmDate{start}至\CiesmDate{end}，按季度推进，提前对接接口与撰写论文。

\paragraph{1—3 月：核查资料并选择参考场}
核查变量、频率并处理资料，固定地形与时间划分、插值及统计校正基线；按研究方法完成\textbf{目标场信息含量检验}与选源。获取与 CMFD 生成站点隔离的国家级台站日值，明确实验室样例、接口与应用方式。交付资料清单、样本索引、粗化算子及信息增量、数据差异报告；频率不足时调整获取计划并记录影响。

\paragraph{4—6 月：检验联合与时间信息的贡献}
完成等信息独立与联合模型，比较地形及单日、邻日输入，固定指标和参数；完成首轮受控实验、预算记录与论文初稿，按风、降水训练稳定性选定专用输出，异常先诊断。

\paragraph{7—9 月：完成约束、台站独立验证与投稿}
完成投影与可行性、软约束强度、地形留出及极端分析，核心稳定且验证支持时检验随机残差校准；完成研究方法所列四组场的台站独立验证，形成结果及失效边界报告，\textbf{并在 9 月底前完成投稿}，预留审稿修改时间。

\paragraph{10—12 月：完成集成应用与成果整理}
用固定模型处理 CIESM 文件，完成格式兼容、实验室联调与确认任务的实际应用；交付规定产品及说明，及时修回论文，汇总程序、配置、数据说明、论文证明与应用记录，提交结题总结。

\subsubsection{预期研究结果}

交付 TerraSR-XJ 模块、可复现处理与训练评估流程及新疆 0.1$^{\circ}$ 逐日四变量试验数据。规格见研究目标，附来源、分辨率含义、范围、误差、台站验证及限制，标明为模型推断结果。

考核目标为研究期内在所属领域重要学术期刊接收或发表论文不少于 1 篇，以及 CIESM 集成与实际应用；以期刊记录、实验室确认的接口运行和任务应用记录验收，获 CIESM 正式接收后由实验室出具应用证明。代码、参数、评价及数据获取说明遵守上游许可和实验室约定，未经许可资料不公开。
